\section{Technical Implementation Details}
\label{app:tech}

\subsection{Checkpoint Structural Index}
\label{app:structural-index}

Before fuzzing a checkpoint \(M_i\), \system builds and freezes
\[
\mathcal{I}(M_i)=(\mathcal{F}_i,\mathcal{Q}_i,
\mathcal{P}_i,\mathcal{Z}_i).
\]
The provenance registry \(\mathcal{P}_i\) contains every stable benchmark
source unit used to construct the checkpoint. A canonical fact in
\(\mathcal{F}_i\) has the form
\[
f=(e,r,v,\tau,p),
\]
where \(p\in\mathcal{P}_i\). The index records the source support and
certification checks for each structured field. Backend-generated identifiers
are stored only in adapter-local mappings and never replace \(p\).

\paragraph{Partial-but-certified extraction.}
Structured fields may come from benchmark schemas, deterministic parsing,
rules, or LLM-assisted extraction. An LLM assertion alone is not a
certificate. Before a field is used for mutation applicability, \system
checks the corresponding benchmark source, span, type, relation support, or
temporal order independently of mutant execution. A field that cannot be
certified remains unstructured. The index may therefore be incomplete without
fabricating metadata.

For a query that admits structured mutation, \(\mathcal{Q}_i\) stores
\[
\operatorname{QI}(q)=(e,s_e,r,s_r,\tau,c,\alpha),
\]
where the spans \(s_e,s_r\) locate the requested entity and relation, \(c\)
stores answer-affecting constraints, and \(\alpha\) records
checkpoint-relative answerability. QueryIntent does not store the expected
answer.

\paragraph{Checkpoint-relative absence.}
An unsupported opportunity requires a certificate
\[
\operatorname{Absent}(z,r,\tau,M_i)=1,
\]
or the corresponding form without \(\tau\). The verification procedure
examines the full checkpoint provenance relevant to the requested entity,
relation, and temporal scope. It does not infer absence merely because no
matching tuple appears in the partial \(\mathcal{F}_i\). The verifier may
combine deterministic matching, structured lookup, and LLM-assisted analysis,
but an LLM assertion alone cannot establish absence. The certificate is fixed
before execution and means only that the request is unsupported by the current
checkpoint.

\paragraph{Structural and evaluator information.}
The structural index contains non-oracle information required for mutation
applicability and canonical region mapping. The preferred flow is
\[
\mathcal{I}(M_i)\longrightarrow
\text{obligation builder}\longrightarrow
\Omega_{r,b}\longrightarrow
\text{scheduler}.
\]
The scheduler normally consumes a frozen projection of
\(\Omega_{r,b}\); this interface choice is not a claim that the structural
index is evaluator-sensitive. The separate ReferenceVault stores gold answers,
\(E^{+}\), \(E^{-}\), correctness annotations, evaluator predicates, and
failure verdicts. Those values are unavailable to mutation construction and
search.

\subsection{Mutation Implementation}
\label{app:mutation-implementation}

Given \(x=(M,q)\) and applicable obligation \(\omega=(x,o,t)\), \system
constructs a candidate according to the fixed operator and target. The target
comes from the frozen structural index; it is not selected using retrieval or
response behavior. An LLM is used only when a query operation requires
free-form realization. It does not choose the operator, target, or memory-state
transition.

\subsubsection{Meaning-Preserving Query Mutation}

Meaning-preserving operations change the surface form while preserving the
entity, relation, temporal scope, answerability, and answer-affecting
constraints in \(\operatorname{QI}(q)\). We use four transformations:
paraphrasing, shortening removable context, adding nonrestrictive context, and
changing syntactic form.

For paraphrasing, the fixed QueryIntent fields are protected while an LLM
expresses the request in different wording. For shortening, \system removes a
known removable span directly or asks an LLM for a shorter realization while
protecting the fixed fields. Nonrestrictive context is accepted only when it
introduces no new constraint. Syntactic-form changes use deterministic
templates when possible and otherwise use an LLM only for grammatical
realization.

\subsubsection{Target-Changing Query Mutation}

A target-changing operation designates one primary field
\(z\in\{e,r,\tau\}\) and a fixed replacement \(z'\). The replacement must be
type-compatible and backed by a certified canonical fact for the unchanged
request structure. The other primary fields and constraints remain fixed.

When the target has a replaceable span, \system applies
\[
q'=q[z\leftarrow z'].
\]
Otherwise, an LLM receives the original QueryIntent and the fixed replacement
only to produce a grammatical form. The scheduler-facing obligation does not
contain the evaluator's expected answer for the replacement. The ReferenceVault
constructs \(G(x')\) separately.

\subsubsection{Unsupported Query Mutation}

An unsupported mutation uses a type-compatible, structurally well-formed
replacement for which
\(operatorname{Absent}(z',r,\tau,M_i)=1\). The target and certificate are
fixed before query realization. Direct slot replacement is used when possible;
otherwise, an LLM may express the already fixed target grammatically.

Neither absence from the partial fact table, failed retrieval, an incorrect
answer, nor an uncorroborated LLM claim establishes unsupportedness. The
evaluator separately assigns the resulting checkpoint-relative reference with
\(a'=\bot\).

\subsubsection{Update Memory Mutation}

An update starts from a certified fact
\[
(e,r,v_{\mathrm{old}},\tau,p)
\]
and fixes a different, type-compatible value \(v_{\mathrm{new}}\) under a
later temporal or state position. The intended transition preserves \(e\) and
\(r\):
\[
(e,r,v_{\mathrm{old}})\longrightarrow(e,r,v_{\mathrm{new}}).
\]
Both values belong to the same canonical region in the original checkpoint
universe \(U_i\).

If the backend exposes a public update operation, \system invokes it. If the
backend represents updates through later interactions or episodes, \system
appends the fixed later interaction and lets the backend process it normally.
It does not directly rewrite backend state into a desired final state. A
backend may retain the old value as history; validity depends on whether its
observable state represents the new value according to its native semantics.

\subsubsection{Deletion Memory Mutation}

A deletion targets an existing certified fact and a faithful public or native
backend operation. Its obligation is applicable only if that operation can
realize the designated transition without changing other information relevant
to the fixed query or otherwise violating the one-input mutation relation.

The adapter may use a backend-local identifier to invoke the operation, but it
must first resolve that identifier to stable checkpoint provenance. Low-level
storage operations are not used to imitate a semantic deletion that the
backend does not support. In particular, Graphiti episode deletion is
applicable only when deleting the episode does not remove multiple relevant
facts or otherwise violate the intended transition; low-level edge deletion is
not substituted. An unsupported deletion opportunity is marked
\textsc{inapplicable}.

Deletion may make the target inactive in the descendant state. Its canonical
region nevertheless remains part of the frozen measurement universe \(U_i\).

\subsubsection{Unrelated Change Memory Mutation}

An unrelated change selects an already indexed region
\[
u\in U_i\setminus
\operatorname{RelevantRegions}(\operatorname{QI}(q)).
\]
It changes the value or state within \(u\) using the backend's normal
operation. All indexed information structurally relevant to the fixed
QueryIntent must remain unchanged. For example, a query about Alice's
residence may be paired with a change from Bob's hobby of tennis to golf when
Bob's hobby is already an indexed region.

The operator cannot introduce a wholly new entity--relation region. A proposed
change to a region outside \(U_i\) is \textsc{inapplicable}. Unrelatedness is
established from QueryIntent and structural regions, not from \(E^{+}\) or
\(E^{-}\).

\subsection{Mutation Applicability and Validation}
\label{app:mutation-validation}

For backend \(b\), an opportunity enters
\(\mathcal{T}_{o,b}(x)\) only if its structural precondition holds and the
backend can faithfully realize it. The resulting applicability mask is frozen
and shared by Random Mutation, Unguided LLM Mutation, and \system for the same
checkpoint, backend, and repetition condition. Exclusion reasons are recorded.

\textsc{Inapplicable} describes an operator-target pair that cannot be
faithfully constructed or executed. Such a target is excluded from
\(\Omega_{r,b}\), is not sent to a generator, and does not enter the applicable
obligation denominator. \textsc{Invalid} describes a generated candidate for
an applicable obligation that fails validation. Invalid generation does not
consume the valid-execution budget, leaves the obligation uncovered, and may
be retried under the cap frozen after the pilot.

For candidate \(x'\) generated from \(x\) by \(o\), validation produces
\[
V_o(x,x')\in\{0,1\}.
\]
Execution occurs only if \(V_o(x,x')=1\). The certificate uses the parent and
candidate forms, structural index, provenance, and observable backend
operation or state information. It never uses the mutant retrieval, mutant
response, ReferenceVault, or failure oracle. An LLM may generate a candidate
or contribute a rejection check, but cannot make the final acceptance decision
alone.

\paragraph{Query validation.}
A meaning-preserving candidate is valid only when its certified QueryIntent
preserves \(e,r,\tau,c,\alpha\). A target-changing candidate is valid only
when the designated slot has the fixed replacement and the remaining primary
slots and constraints are unchanged. Unsupported validation additionally
requires the precomputed checkpoint-wide absence certificate.

\paragraph{Memory-state validation.}
Memory validation checks \(M\rightarrow M'\) independently of query behavior.
An update must preserve entity and relation, use the fixed new value at a later
state position, remain in the target's anchored region, and be observably
represented under native backend semantics. A deletion must show that the
designated target existed before the operation and became removed or inactive
afterward with no unintended relevant change. An unrelated change must remain
within its fixed existing region and leave all structurally relevant indexed
information unchanged.

When the backend cannot provide operation results or observable state
sufficient for these checks, the opportunity is \textsc{inapplicable} for
RQ1. It is not admitted as an unverified behavioral probe.

\subsection{Canonical Region Construction and Mapping}
\label{app:region-mapping}

For a certified fact, define
\[
\kappa(f)=(\texttt{entity-relation},
\operatorname{norm}(e),\operatorname{norm}(r)).
\]
The normalization procedure is chosen after the pilot and frozen before full
evaluation. It is shared across methods.

For every stable checkpoint source unit \(p\), let
\[
K_i(p)=\{\kappa(f):f\text{ has provenance }p
\text{ and a certified entity--relation key}\}.
\]
The source contributes
\[
\operatorname{RegionsForSource}_i(p)=
\begin{cases}
K_i(p), & K_i(p)\neq\varnothing,\\
\{(\texttt{source},p)\}, & K_i(p)=\varnothing.
\end{cases}
\]
Thus, a source with any certified semantic key contributes only its certified
semantic regions. Partial residual text does not cause a simultaneous source
fallback. When the source yields no certified entity--relation region, the
stable source ID is the sole conservative fallback.

The original checkpoint defines one frozen measurement anchor
\[
U_i:=U(M_i)=
\bigcup_{p\in\mathcal{P}_i}\operatorname{RegionsForSource}_i(p).
\]
No descendant state defines another universe. For every retrieved entry \(m\)
from any descendant state, the backend adapter resolves local identifiers to
stable provenance and constructs
\[
\operatorname{FactMap}_{b,i}(m)=
\{f:m\text{ can reliably be shown to represent }f\}.
\]
It then defines
\[
R^{\mathrm{sem}}_{b,i}(m)=
\{\kappa(f):f\in\operatorname{FactMap}_{b,i}(m)\}.
\]
Source fallback is represented separately:
\[
R^{\mathrm{fallback}}_{b,i}(m)=
\{(\texttt{source},p):m\text{ has reliable provenance to }p
\land K_i(p)=\varnothing\}.
\]
The complete mapping is
\[
\rho_{b,i}(m)=R^{\mathrm{sem}}_{b,i}(m)\cup
R^{\mathrm{fallback}}_{b,i}(m)\subseteq U_i.
\]
If source \(p\) has no certified semantic region, reliable provenance to
\(p\) may map \(m\) to \((\texttt{source},p)\). If \(p\) has exactly one
certified semantic region, reliable provenance to \(p\) is sufficient for
that mapping. If \(p\) has multiple semantic regions, \(m\) maps only to the
subset of canonical facts that its content and provenance reliably identify.
When such an entry cannot be disambiguated, the adapter marks it ambiguous or
unmapped for RegionCov instead of assigning all of \(K_i(p)\). A summarized
entry may map to several regions only when it reliably represents several
certified facts. Values created by a valid memory mutation inherit the
canonical key of their designated indexed target. If reliable checkpoint
provenance is absent, the adapter records the entry as unmapped and does not
invent a region. RegionCov therefore counts regions represented in the
retrieved entry rather than every region in its source unit.

\paragraph{Mapping-quality diagnostic.}
For all retrieved entries considered in a condition, we report
\[
\operatorname{MappingRate}=
\frac{|\{m:\rho_{b,i}(m)\neq\varnothing\}|}
{|\{m:m\text{ is a retrieved entry considered}\}|}.
\]
We also log mapped entries, ambiguous entries whose multi-fact semantic
content cannot be reliably disambiguated, and unmapped entries lacking another
reliable mapping. MappingRate is reported by backend/benchmark condition and
optionally by method as a sanity check. It is diagnostic only: it does not
guide scheduling, seed retention, or mutation selection; it does not enter the
RegionCov numerator or denominator; and it is not treated as a primary RQ1
effectiveness metric.

\subsection{Campaign Initialization and Execution Accounting}
\label{app:campaign-accounting}

One campaign is identified by
\[
(M_i,b,h,s),
\]
for checkpoint \(M_i\), backend \(b\), method \(h\), and repetition seed \(s\).
Its budget \(B\) is the number of valid new mutant executions.

For fixed \((M_i,b,s)\), \system first constructs one backend checkpoint and a
frozen initialization artifact. Exact state cloning for all three methods is
preferred. If cloning is unavailable, all methods replay the same artifact
with all controllable randomness, model versions, prompts, decoding settings, and
ingestion configuration frozen. Identical replay input alone does not prove
equivalent backend state. Before method-specific fuzzing, \system checks
initialization equivalence and records a stable state identity, fingerprint,
or canonical state projection sufficient to establish equivalent starting
states. The concrete fingerprinting algorithm is deferred to the backend
implementation and capability-test phase. If equivalence cannot be
established, the condition is recorded as an initialization or capability
failure and excluded from paired RQ1 comparison until resolved.

After equivalence is established, the original query set is executed to
produce shared parent observations and
\[
C_0^{(i,b,s)}.
\]
These executions do not consume \(B\).

An observation cache is keyed by exact state identity and exact query
identity. An initial seed receives its cached \(U(x)\) during shared
initialization. When a valid mutant is executed, its \(U(x')\) is stored with
the retained seed. If that exact seed later becomes a parent, the stored
observation is reused, and only the new mutant execution consumes one unit of
\(B\).

State and query identity must be proven before reuse. If a backend state must
be reconstructed or identity cannot be established, \system logs the replay,
reconstruction, or parent execution separately. Such overhead is neither
silently reused nor silently counted as a valid mutant probe. The same identity
and overhead protocol is used by all methods.

Each campaign stops after exactly \(B\) valid new mutant executions. Invalid
generation and inapplicable opportunities consume no units. For \(N\)
checkpoints, each fixed backend/method/repetition condition therefore performs
\(N B\) valid new mutant executions, plus separately reported initialization
and reconstruction overhead.

\subsection{Manual Audit and Information-Flow Check}
\label{app:audit-information-flow}

After the automatic pipeline operates, we conduct a stratified manual audit
across benchmark, mutation operator, and backend where backend semantics affect
applicability. The audit covers QueryIntent and span extraction, canonical
facts, provenance mapping, target-changing validity, checkpoint-relative
absence, fact-level retrieval-to-region mapping accuracy, update relations,
deletion applicability, and unrelated-change validity. We report the sampling
procedure, sample counts, exclusion reasons, and empirical metadata and
mutation validity rates. We additionally report retrieval-mapping accuracy,
MappingRate, and mapped, ambiguous, and unmapped entry counts separately from
the primary effectiveness metrics. The sample size is chosen after the pilot
and frozen before full evaluation; the audit is quality control rather than
the primary construction mechanism.

We also test the search/evaluator boundary. Given a sealed observable
execution trace, we permute evaluator-only gold answers and failure labels and
replay scheduling and retention from the same retrieval observations. Every
scheduling and retention decision must remain identical. The executable
definitions of \(\mathcal{C}_o\) and answer matching, and the distinct-failure
deduplication rule, are specified separately before full evaluation rather
than inferred by the implementation.
